% ===================================================================
%  அட்டவணை எண் 3 — இளமையும் வாலிபமும்.
%
%  Traduction, faite en 2026, en tamoul d'aujourd'hui, sur l'ido de
%  texto/io. Regles de la colonne : voir 10-attavanai-01.tex. Deux
%  scenes, deux numerotations : les cles portent c1 et c2.
%
%  DIX-SEPT RETOURNEMENTS SUR VINGT-CINQ PAIRES — le meme compte que
%  la colonne coreenne au meme tableau, et pour la meme raison. C'est
%  un DEFILE, et un defile se decrit par attributs : la nourrice AVEC
%  sa coiffe, le mendiant APPUYE sur sa bequille, la vieille APPUYEE
%  sur son baton, le meunier AVEC son chapeau, les abeilles DE la
%  ruche, la fleche DU clocher, le clocher QUI couronne l'eglise.
%  Chacun de ces liens s'inverse en tamoul comme en coreen. Le
%  tableau 2, qui enumerait, en coutait un seul.
%
%  ET J'EN AVAIS COMPTE SEIZE. J'avais range « (batel)-turniras (32)
%  proxim la kayo (33) » parmi les exemptions, en me disant qu'un
%  complement de lieu n'est pas un modificateur de nom et que l'ordre
%  libre du tamoul le laisserait ou il est. C'est faux : le tamoul
%  place le circonstanciel AVANT le verbe et donc avant l'objet, et
%  la phrase est sortie « pres du quai (33), ils tournoient (32) ».
%  renvoji.py l'a relevee au premier passage. La correction tient en
%  une coupure — on nomme le tournoi, puis on dit ou il a lieu — et
%  la lecon est pour le classement, pas pour la phrase : un
%  complement de lieu N'EST PAS une exemption. Il n'y en a toujours
%  que trois : l'enumeration, le renvoi de gauche deja modificateur,
%  la frontiere de phrase.
%
%  ET DEUX LANGUES SANS PARENTE QUI TOMBENT SUR 17/25 AU MEME
%  TABLEAU, CE N'EST PAS UNE COINCIDENCE DE PLUS — c'est la
%  confirmation de ce que le tableau 14 coreen avait fini par dire :
%  la liste de parigi.py mesure l'ido, et les trois motifs
%  d'exemption sont des proprietes de l'ido. Deux langues a tete
%  finale qui lisent la meme liste font le meme partage.
%
%  LE FAC-SIMILE DONNE DEUX FOIS LE RENVOI (18) ET DEUX FOIS LE (19)
%  au bloc c1-09-1 : les arbres fruitiers puis les pechers pour le
%  premier, le presbytere puis les abricotiers pour le second. C'est
%  une faute de l'edition ido, et la source ne bouge pas : la
%  traduction porte les quatre, dans l'ordre recu.
%
%  LE MOT DEVENU UNE INSULTE, PREMIERE FOIS DANS CETTE COLONNE. La
%  baraque du bloc c2-14-1 montre « la giganto e la nano ». Le tamoul
%  a குள்ளன், qui traduit exactement le second et qui sert
%  aujourd'hui d'injure courante. La colonne ecrit மிகச் சிறிய
%  மனிதர், et — pour que la paire reste une paire — மிகப் பெரிய
%  மனிதர் pour le geant, dont les mots disponibles (அரக்கன்,
%  இராட்சதன்) designent des demons et non des hommes. Meme
%  arbitrage que pour le 난쟁이 coreen et le కసాయి telougou. La
%  baraque, elle, reste : c'est la CHOSE.
%
%  LES SABOTS (75) SONT LE TROISIEME PIEGE DE CETTE COLONNE, apres le
%  jeu de barres et les quilles du tableau 2. Le tamoul a பாதுகை, la
%  sandale de bois, objet reel et vivant — et ce n'est pas un sabot :
%  c'est une semelle a tenon, pas une chaussure creusee. La colonne
%  ecrit மரக் காலணிகள். Troisieme fois en trois tableaux qu'un mot
%  juste designe un autre objet, et troisieme refus.
%
%  CE QUI S'EMPRUNTE ICI S'EMPRUNTE PARCE QUE LA CHOSE NE POUSSE PAS
%  AU PAYS TAMOUL : பாப்ளர் (3) le peuplier, பிளாட்டன் (17) le
%  platane, பீச் (18) le pecher, ஆப்பிரிகாட் (19) l'abricotier,
%  துலிப் (36) la tulipe, கார்னேஷன் (34) l'oeillet. La coupure suit
%  la carte, comme au tableau 9 coreen et au tableau 9 telougou.
%
%  ET DEUX NOMS SONT ECRITS EN SACHANT CE QU'ILS COUTENT : le muguet
%  (37), qui n'a pas de nom tamoul, est ecrit மே மலர் — la fleur de
%  mai, ce que l'ido lui-meme a fait en ecrivant « maiflori » ; et
%  l'hirondelle (4) est ecrite தகைவிலான், qui est le nom tamoul de
%  l'oiseau mais un nom de livre. Le telougou avait du renoncer a la
%  sienne. Un nom rare vaut mieux qu'un nom faux ; un nom invente ne
%  vaut rien.
%
%  LE VOCABULAIRE CATHOLIQUE EXISTE EN TAMOUL DEPUIS QUATRE SIECLES
%  et la colonne s'en sert : பங்குத் தந்தை (15) le cure, பங்கு
%  இல்லம் (19) le presbytere, ஞானத் தந்தை (49) le parrain, ஞானத்
%  தாய் (53) la marraine, ஞானஸ்நானம் le bapteme. Le Tamil Nadu a des
%  paroisses reelles, et ce sont leurs mots — pas des calques faits
%  pour l'occasion.
%
%  ET « julio » EST EN GRAS AU BLOC c2-01-1, comme dans les trente-
%  huit colonnes qui precedent. La ligne « ta » a ete posee dans
%  gravuri/korekti.json AVANT que ce tableau soit ecrit, et non six
%  colonnes plus tard.
% ===================================================================

%%K t03-apar-1 apar
\VUsaut{3.0mm}
\VUtitre{15.0pt}{60}{அட்டவணை எண் 3}
\VUsaut{1.6mm}
\VUtitre{11.4pt}{0}{இளமையும் வாலிபமும்.}
\VUsaut{3.4mm}

%%K t03-apar-2 apar
(3 - ஆம், 4 - ஆம் அட்டவணைகள், நான்கு \VUgras{குடும்பக் காட்சிகளில்}
\VUgras{வானிலை}, \VUgras{வயது}, \VUgras{குடும்பம்}, \VUgras{உடையணிதல்},
\VUgras{நிலைமை} பற்றிய அடிப்படைக் கருத்துகளைக் காட்டும் வகையில் ஒன்றோடொன்று
இணைக்கப்பட்டிருக்கின்றன.)

\VUsaut{4.0mm}
%%K t03-tit-1 sub
\VUcentre{12.0pt}{0}{\textit{முதல் காட்சி.}}
\VUsaut{3.0mm}

%%K t03-tit-2 sub
\VUpk{10.2pt}{40}{கிராமத்தில் ஞானஸ்நான வைபவம்.}
\VUsaut{2.4mm}

%%K t03-tit-3 sub
\VUpk{10.2pt}{40}{வசந்த காலம்.}
\VUsaut{3.0mm}

%%K t03-c1-01-1 p
1. --- நாம் \VUgras{வசந்த காலத்தில்} இருக்கிறோம் ; \VUgras{மார்ச்},
\VUgras{ஏப்ரல்} அல்லது \VUgras{மே} மாதங்களில், \VUgras{வாரத்தின்} நாள்களான
\VUgras{திங்கள்}, \VUgras{செவ்வாய்} அல்லது \VUgras{புதன்} கிழமைகளில். காலை
\VUgras{பத்து மணி}.

%%K t03-c1-02-1 p
2. --- \VUgras{வானிலை} அழகாக இருக்கிறது, \VUgras{காற்று} தூய்மையாக இருக்கிறது.
சூரியன் ஒளிர்கிறது. சில \VUgras{மேகத் துணுக்குகள்} \textsuperscript{(2)}
நீல \VUgras{வானத்தில்} \textsuperscript{(1)} மேலெழுகின்றன.

%%K t03-c1-03-1 p
3. --- \VUgras{மரங்களுக்கு} \VUgras{இலைகள்} இன்னும் கிட்டத்தட்ட இல்லை. மெல்லிய
\VUgras{பாப்ளர் மரங்களுக்கும்} \textsuperscript{(3)} உயரமான \VUgras{பிளாட்டன்
மரத்திற்கும்} \textsuperscript{(17)} இதுவரை அரும்புகள் மட்டுமே உண்டு.

%%K t03-c1-04-1 p
4. --- \VUgras{தகைவிலான்கள்} \textsuperscript{(4)} திரும்பி வந்து மகிழ்ச்சியான
கீச்சொலியுடன் \VUgras{கூரிய முனையைச்} \textsuperscript{(7)} சுற்றிப்
பறக்கின்றன ; அந்த முனை \VUgras{மணிக் கோபுரத்தின்} \textsuperscript{(5)} உச்சியில்
இருக்கிறது, அந்த மணிக் கோபுரம் கிராமத்தின் \VUgras{தேவாலயத்திற்கு}
\textsuperscript{(6)} முடி சூட்டுகிறது.

%%K t03-tit-4 sub
\VUsaut{3.4mm}
\VUpk{10.2pt}{40}{தேவாலயம்.}
\VUsaut{3.0mm}

%%K t03-c1-05-1 p
5. --- அந்தத் தேவாலயம் மிகவும் எளிமையானது ; அதற்குப் பெரிய
\VUgras{பெருங்கதவும்} \textsuperscript{(13)} பருமனான \VUgras{கடிகாரமும்}
\textsuperscript{(8)} உண்டு. சிறிய \VUgras{முள்} \textsuperscript{(9)} X என்னும்
எண்ணின் மேல் இருக்கிறது ; அது \VUgras{மணிகளைக்} காட்டுகிறது.
\VUgras{பெரிய முள்} \textsuperscript{(10)} XII இன் மேல் இருக்கிறது (நண்பகல்) ;
அது \VUgras{நிமிடங்களைக்} காட்டுகிறது.

%%K t03-c1-06-1 p
6. --- மணிக் கோபுரத்தில் \VUgras{மணிகள்} \textsuperscript{(11)} மகிழ்ச்சியாக
ஒலிக்கின்றன. கூரையின் உச்சியில் கல்லால் ஆன சிறிய \VUgras{சிலுவை}
\textsuperscript{(12)} நிற்கிறது.

%%K t03-c1-07-1 p
7. --- தேவாலயத்திற்குப் பெரிய \VUgras{பெருங்கதவு} \textsuperscript{(13)}
மட்டும் அல்ல, பக்கத்தில் ஒரு சிறிய கதவும் உண்டு. அந்தச் சிறிய கதவின் வழியாக,
தன் \VUgras{நீளங்கியை} அணிந்த \VUgras{பங்குத் தந்தை} \textsuperscript{(15)}
\VUgras{பூசைப் பொருள் அறையிலிருந்து} \textsuperscript{(14)} வெளியே வந்து
\VUgras{பங்கு இல்லத்திற்குப்} \textsuperscript{(19)} போகிறார்.

%%K t03-c1-08-1 p
8. --- அவருக்கு அருகில், இதோ \VUgras{பால்காரி} \textsuperscript{(24)} ; அவளுடன்
அவளுடைய \VUgras{கழுதை} \textsuperscript{(23)}. அவள் நகரத்திலிருந்து திரும்பி
வருகிறாள் ; அங்கே குழந்தைகளுக்கு நல்ல பால் கொண்டு போயிருந்தாள்.

%%K t03-tit-5 sub
\VUsaut{4.0mm}
\VUpk{10.2pt}{40}{பழத் தோட்டம்.}
\VUsaut{3.4mm}

%%K t03-c1-09-1 p
9. --- அலிபோரோன் ஐயா (கழுதை) அந்தக் காலை நடையைப் பற்றி முற்றிலும் மனநிறைவாக
இருக்கிறார். அவர் மணம் கமழும் காற்றை இன்பத்துடன் இழுத்து விடுகிறார் ; அந்த
மணம் \VUgras{மலர்களுடையது} \textsuperscript{(20)}, அவை பக்கத்துப்
\VUgras{பழத் தோட்டத்தின்} \VUgras{பழ மரங்களை} \textsuperscript{(18)}
மூடியிருக்கின்றன. அந்த \VUgras{பீச் மரங்களும்} \textsuperscript{(18)}
\VUgras{ஆப்பிரிகாட் மரங்களும்} \textsuperscript{(19)} முற்றிலும் ரோஜா
நிறமாகவும் வெண்மையாகவும் இருக்கின்றன ; நல்ல அறுவடையை அவை நம்ப வைக்கின்றன.

%%K t03-c1-10-1 p
10. --- அதே வேளையில் சிறிய \VUgras{தேனீக்கள்} \textsuperscript{(22)} அங்கே
கொள்ளையடிக்கப் போகின்றன ; அவை \VUgras{தேன் கூட்டைச்} \textsuperscript{(21)}
சேர்ந்தவை, ரீங்காரமிட்டுத் தங்கள் இனிய \VUgras{தேனைச்} சேர்க்கின்றன.

%%K t03-c1-11-1 p
11. --- ஆனால் என்ன நடக்கிறது? இதோ கிராமத்துச் சிறுவர்கள் ஓடி வந்து பயந்த
\VUgras{வாத்துகளை} \textsuperscript{(25)} விரட்டுகிறார்கள். நோட்டரியின்
மகனுக்கு \VUgras{ஞானஸ்நானம்} முடிந்த பிறகு கூட்டம் தேவாலயத்திலிருந்து
வெளியே வருவதுதான் அது.

%%K t03-c1-12-1 p
12. --- சிறிய யாகோபுஸ் தன் \VUgras{தொட்டிலில்} \textsuperscript{(26)}
மணிகளின் மகிழ்ச்சியான ஒலியால் விழித்துக் கொள்கிறான் ; அவனுடைய தமக்கை
மக்தலேனா, அவளும் ஞானஸ்நான ஊர்வலத்தைப் பார்க்க விரும்பி, தன் தலைமுடியை
வாரவும் தன் உடையை அணிவிக்கவும் வீட்டு வாசற்படிக்கு வரும்படி தன் தாயை
வேண்டுகிறாள்.

%%K t03-c1-13-1 p
13. --- தாய் ஒப்புக் கொள்கிறாள். அவள் தன் \VUgras{முக்காட்டையும்}
\textsuperscript{(28)}, தன் \VUgras{ரவிக்கையையும்} \textsuperscript{(29)}, தன்
\VUgras{முன்னாடையையும்} \textsuperscript{(30)} அணிந்து கொண்டு, கதவின் முன்
சிறிய நாற்காலியில் வந்து அமர்கிறாள். மக்தலேனாவிடம் தன் \VUgras{உள்பாவாடையும்}
\textsuperscript{(31)} தன் \VUgras{இறுக்குக் கச்சையும்} \textsuperscript{(32)}
மட்டுமே இருக்கின்றன.

%%K t03-c1-14-1 p
14. --- அவள் எவ்வளவு மகிழ்ச்சியாக இருக்கிறாள் ! தான் விரும்பும் மலர்களை
மறந்து விடுகிறாள் — தன் \VUgras{கார்னேஷன்களை} \textsuperscript{(34)}, அவற்றின்
மேல் \VUgras{வண்ணத்துப் பூச்சிகள்} \textsuperscript{(35)} வந்து அமர்கின்றன ;
தன் \VUgras{துலிப் மலர்களை} \textsuperscript{(36)} ; தன் \VUgras{மே மலர்களை}
\textsuperscript{(37)} — \VUgras{தொட்டிகளில்} \textsuperscript{(38)},
\VUgras{சுவரின்} \textsuperscript{(33)} மேல் — ; தன் \VUgras{ரோஜாக்களை}
\textsuperscript{(39)}, அவை மிகவும் அழகான வேலியை உண்டாக்குகின்றன. ஊர்வலத்தில்
வரும் பெண்களின் செல்வமிக்க உடைகளை அவள் கவனிக்கிறாள்.

%%K t03-c1-15-1 p
15. --- கிராமத்துச் சிறுவர்கள் உடைகளைப் பெரிதாகக் கவனிப்பதில்லை.
\VUgras{மிட்டாய்களையோ} \textsuperscript{(55)} \VUgras{சில்லறை நாணயங்களையோ}
\textsuperscript{(52)} அடைவதற்காக ஒருவரையொருவர் தள்ளிக் கொண்டு விரைகிறார்கள்.
அவர்களுள் ஒருவன் அழுகிறான் \textsuperscript{(40)}.

%%K t03-c1-16-1 p
16. --- இன்னொருவன் \VUgras{நாயின்} \textsuperscript{(42)} காலை மிதித்து
விட்டான் ; அது கத்திக் கொண்டு ஓடி, \VUgras{கோழியையும்} \textsuperscript{(43)}
அதன் \VUgras{குஞ்சுகளையும்} \textsuperscript{(44)} விரட்டுகிறது.

%%K t03-tit-6 sub
\VUsaut{5.0mm}
\VUpk{10.2pt}{40}{ஞானஸ்நான ஊர்வலம்.}
\VUsaut{4.0mm}

%%K t03-c1-17-1 p
17. --- ஊர்வலத்தின் முன்னால் \VUgras{செவிலித் தாய்} \textsuperscript{(45)}
நடக்கிறாள். அவளிடம் நீண்ட நாடாக்கள் கட்டிய அழகான \VUgras{தலை உறை}
\textsuperscript{(46)} இருக்கிறது. அவள் \VUgras{பிறந்த குழந்தையைச்}
\textsuperscript{(47)} சுமக்கிறாள் ; அந்தக் குழந்தை முழுவதும்
\VUgras{நூற்பின்னலால்} \textsuperscript{(48)} உடுத்தப்பட்டிருக்கிறது.

%%K t03-c1-18-1 p
18. --- செவிலித் தாய்க்குப் பின்னால் \VUgras{ஞானத் தந்தையும்}
\textsuperscript{(49)} \VUgras{ஞானத் தாயும்} \textsuperscript{(53)}
வருகிறார்கள். அவர்கள் மிட்டாய்களையும் சில்லறை நாணயங்களையும்
பங்கிடுகிறார்கள். ஞானத் தந்தையிடம் \VUgras{கையுறைகளும்} \textsuperscript{(51)}
\VUgras{உயரத் தொப்பியும்} \textsuperscript{(50)} இருக்கின்றன. ஞானத் தாய்
கையில் அழகான \VUgras{பூங்கொத்தைப்} \textsuperscript{(54)} பிடித்திருக்கிறாள்.

%%K t03-c1-19-1 p
19. --- அவர்களுக்குப் பின்னால் குழந்தையின் \VUgras{மாமாவையும்}
\textsuperscript{(56)} \VUgras{அத்தையையும்} \textsuperscript{(58)}
பார்க்கிறோம். அத்தையிடம் நேர்த்தியான \VUgras{தொப்பி} \textsuperscript{(59)}
இருக்கிறது ; மாமாவிடம் கருப்பு \VUgras{நீள மேலங்கியும்} \textsuperscript{(57)}
வெள்ளை \VUgras{உள்ளங்கியும்} இருக்கின்றன. இதோ பிறகு \VUgras{ஒன்றுவிட்ட
சகோதரனும்} \textsuperscript{(60)} \VUgras{ஒன்றுவிட்ட சகோதரியும்}
\textsuperscript{(61)}. இவள் பிறந்த குழந்தையின் \VUgras{சகோதரியை}
\textsuperscript{(62)}, சிறிய பெர்த்தாவை, கையால் பிடித்திருக்கிறாள்.

%%K t03-c1-20-1 p
20. --- பெர்த்தாவின் இன்னொரு \VUgras{சகோதரன்} \textsuperscript{(63)}
பின்னால் நடக்கிறான். அவன் ஒரு \VUgras{உறவினருக்குக்} \textsuperscript{(64)}
கை கொடுக்கிறான். குடும்பத்தின் \VUgras{நண்பர்கள்} \textsuperscript{(65)}
பிறகு வருகிறார்கள்.

%%K t03-tit-7 sub
\VUsaut{4.6mm}
\VUpk{10.2pt}{40}{பார்ப்பவர்கள்.}
\VUsaut{3.4mm}

%%K t03-c1-21-1 p
21. --- ஊர்வலத்திற்கு அருகில், இதோ ஒரு \VUgras{பிச்சைக்காரர்}
\textsuperscript{(66)} ; அவர் தன் \VUgras{ஊன்றுகோலின்} \textsuperscript{(67)}
மேல் சாய்ந்திருக்கிறார். அவர் தர்மம் கேட்கிறார் (பிச்சை எடுக்கிறார்).

%%K t03-c1-22-1 p
22. --- மறுபக்கத்தில் மிகவும் \VUgras{பணிவான} \textsuperscript{(68)} ஒரு
சிறுவன் இருக்கிறான். அவனுக்குத் தெரிந்தவர்களை வணங்கி « \VUgras{வணக்கம்} »
என்று வாழ்த்துகிறான்.

%%K t03-c1-23-1 p
23. --- ஞானஸ்நான ஊர்வலத்தைப் பார்ப்பதற்காகக் கிராமத்தார் தங்கள்
வீடுகளிலிருந்து வெளியே வந்திருக்கிறார்கள். ஒரு \VUgras{கிராமத்தானைப்}
\textsuperscript{(69)} பார்க்கிறோம் ; அவனிடம் \VUgras{பருத்திக் குல்லாய்}
இருக்கிறது, அவனுக்குப் பக்கத்தில் ஒரு \VUgras{கிராமத்துப் பெண்}
\textsuperscript{(70)}. பிறகு ஒரு \VUgras{கிழவி} \textsuperscript{(71)} ; அவள்
தன் \VUgras{கைத்தடியின்} \textsuperscript{(72)} மேல் சாய்ந்திருக்கிறாள்.
கடைசியாக \VUgras{மாவு ஆலைக்காரர்} \textsuperscript{(73)} ; அவரிடம் பெரிய
\VUgras{கம்பளித் தொப்பி} \textsuperscript{(74)} இருக்கிறது. \VUgras{மரக்
காலணிகள்} \textsuperscript{(75)} அணிந்த ஒரு இளம் கிராமத்துப் பெண்ணும்
இருக்கிறாள் ; அவள் தன் குழந்தைக்கு நடக்கக் கற்றுத் தருகிறாள்.

%%K t03-c1-24-1 p
24. --- இந்தக் காட்சி மிகவும் மகிழ்ச்சி தருவது.

\VUsaut{4.0mm}
%%K t03-tit-8 sub
\VUcentre{12.0pt}{0}{\textit{இரண்டாம் காட்சி.}}
\VUsaut{3.0mm}

%%K t03-tit-9 sub
\VUpk{10.2pt}{40}{பொது விழா.}
\VUsaut{2.4mm}

%%K t03-tit-10 sub
\VUpk{10.2pt}{40}{நீர் விளையாட்டுகளும் படகுப் போட்டிகளும்.}
\VUsaut{3.0mm}

%%K t03-c2-01-1 p
1. --- \VUgras{கோடை} வந்து விட்டது. \VUgras{ஜூன்} மாதத்தின் (ஜூலை அல்லது
\VUgras{ஆகஸ்ட்}) கடும் வெயில் இளைஞர்களைக் குளிக்கவும் படகு
செலுத்தவும் தூண்டுகிறது.

%%K t03-c2-02-1 p
2. --- ஆற்றின் வலப் புறக் கரையில், படகுக்காரர்கள் நீர் விளையாட்டுகளிலும்
படகுப் போட்டிகளிலும் கலந்து கொள்வதற்காக உடைகளைக் களைகிறார்கள்.

%%K t03-c2-03-1 p
3. --- அவர்களுள் ஒருவன் தன் \VUgras{காலணிகளைக்} \textsuperscript{(1)}
கழற்றுகிறான் ; அவற்றின் கயிற்றை அவிழ்க்கிறான் (பொத்தானைக் கழற்றுகிறான்).
அவனுடைய இரண்டு தோழர்கள் ஏற்கெனவே தயாராக இருக்கிறார்கள். இப்போது அவர்களிடம்
தங்கள் \VUgras{பின்னலாடையும்} \textsuperscript{(2)} \VUgras{நீச்சல் உடையும்}
\textsuperscript{(3)} மட்டுமே இருக்கின்றன. தங்கள் நண்பர்களுக்குக் காத்திருந்து
பேசிக் கொள்கிறார்கள். விடுமுறையைப் பற்றியும் மறு கரையில் நடக்கும் விழாவைப்
பற்றியும் பேசுகிறார்கள்.

%%K t03-c2-04-1 p
4. --- அவர்களுக்கு அருகில் தரையில் அவர்களுடைய தோழர்களின் \VUgras{உடைகள்}
கிடக்கின்றன : ஒரு ஜோடி \VUgras{குறுங் காலணிகள்} \textsuperscript{(4)},
\VUgras{செருப்புகள்} \textsuperscript{(5)}, \VUgras{காலுறைகள்}
\textsuperscript{(6)}, \VUgras{கம்பளி} \textsuperscript{(7)},
\VUgras{வைக்கோல்} \textsuperscript{(8)} தொப்பிகள், \VUgras{கழுத்துப் பட்டை}
\textsuperscript{(9)}.

%%K t03-c2-05-1 p
5. --- இளைஞர்களுள் ஒருவன் தன் \VUgras{மேலுடையைக்} \textsuperscript{(10)}
கவனமாக மடித்தான். அதை ஒரு துவைத் துணியின் மேல் வைக்கிறான் ; அதற்குள்
அவனுக்கு அடுத்திருப்பவன் விரைவில் இரண்டு \VUgras{உடைகளையும்} வைப்பான்.

%%K t03-c2-06-1 p
6. --- ஆறாவது சிறுவன் தாமதிக்கிறான். அவனுடைய \VUgras{குல்லாய்}
\textsuperscript{(11)} இன்னும் தலையில் இருக்கிறது. அவன் தன் \VUgras{சட்டையையோ}
\textsuperscript{(12)}, தன் \VUgras{சட்டைக் கழுத்தையோ} \textsuperscript{(13)},
தன் \VUgras{கைப் பட்டைகளையோ} \textsuperscript{(14)} கழற்றவில்லை. கையில் தன்
\VUgras{உள்ளங்கியைப்} \textsuperscript{(17)} பிடித்திருக்கிறான் ; அவனுடைய
\VUgras{கால்சட்டையும்} \textsuperscript{(16)} \VUgras{தோள்வார்களும்}
\textsuperscript{(15)} இன்னும் இருக்கின்றன. அடுத்த ஓட்டத்திற்கு அவன்
தயாராக இருக்க மாட்டான், நிச்சயம்.

%%K t03-c2-07-1 p
7. --- இளைஞர்களுக்கு அருகில் ஒரு சிறு பாதை ஆற்றின் விளிம்பிற்கு
அழைத்துச் செல்கிறது ; அதன் இரு பக்கங்களிலும் \VUgras{முள்ளுச் செடிகள்}
\textsuperscript{(18)} நிறைய இருக்கின்றன. அங்கே யாகோபுஸ் அப்பா
\VUgras{நாணல்களுக்கு} \textsuperscript{(19)} இடையில் \VUgras{வானவேடிக்கையைத்}
\textsuperscript{(20)} தயார் செய்கிறார் ; மாலையில் அதைக் கொளுத்துவார்கள்.

%%K t03-tit-11 sub
\VUsaut{4.4mm}
\VUpk{10.2pt}{40}{ஓட்டம்.}
\VUsaut{3.4mm}

%%K t03-c2-08-1 p
8. --- ஆற்றின் மேல் சில பெண்கள் தங்கள் குடைகளால் தங்களை மறைத்துக் கொண்டு
\VUgras{படகில்} \textsuperscript{(21)} உலா வருகிறார்கள். மிகவும் சுவாரசியமான
ஓட்டத்தை அவர்கள் பின்தொடர்கிறார்கள். ஒவ்வொரு \VUgras{சிறு படகிலும்}
\textsuperscript{(22)} நான்கு \VUgras{துடுப்பு வீரர்கள்} \textsuperscript{(24)}
தங்கள் முழு வலிமையுடன் துடுப்பு போடுகிறார்கள் ; அதே வேளையில்
\VUgras{சுக்கான் பிடிப்பவர்} \textsuperscript{(26)} \VUgras{சுக்கானைப்}
\textsuperscript{(25)} பிடித்து அந்த மெல்லிய படகைச் செலுத்துகிறார்.
\VUgras{துடுப்புகள்} \textsuperscript{(23)} தாளத்துடன் நீரை அடிக்கின்றன ;
இலேசான படகுகள் தலை சுற்றும் வேகத்துடன் செல்கின்றன.

%%K t03-c2-09-1 p
9. --- சில இளைஞர்கள் பாலத்தின் தூணுக்கு அருகில் \VUgras{படகு முன் கழியின்}
\textsuperscript{(28)} பரிசுக்காகப் போட்டியிடுகிறார்கள். அவர்களுள் ஒருவன்
நுனியில் கட்டப்பட்ட சிறிய \VUgras{கொடியை} \textsuperscript{(29)} அடைவதற்காக,
வளையும் வழுக்கும் அந்தக் கழியின் மேல் எச்சரிக்கையாக நடக்கிறான். அதிர்ஷ்டம்
இல்லாத அவனுடைய \VUgras{போட்டியாளர்} \textsuperscript{(30)} நீரில் விழுந்து
விட்டான். கரை சேர அவன் நீந்துகிறான்.

%%K t03-c2-10-1 p
10. --- வயது கூடிய இளைஞர்கள் \VUgras{படகுச் சண்டை} \textsuperscript{(32)}
நடத்துகிறார்கள் ; அது \VUgras{கற்கரைக்கு} \textsuperscript{(33)} அருகில்
நடக்கிறது. இந்த
விளையாட்டுகள் எல்லாவற்றையும் ஏராளமான \VUgras{மக்கள் கூட்டம்}
\textsuperscript{(31)} பார்க்கிறது ; அது \VUgras{பாலத்தின்}
\textsuperscript{(27)} தடுப்புச் சுவரின் மேல் சாய்ந்தும் \VUgras{கரையில்}
திரண்டும் இருக்கிறது. இருந்தாலும் சில பார்வையாளர்கள் \VUgras{பரிசுக்
கழியின்} \textsuperscript{(45)} விளையாட்டைப் பின்தொடரவோ, \VUgras{பரிசு
வழங்குவதற்காக} வரும் \VUgras{மேயர் ஊர்வலத்தைப்} போற்றவோ திரும்பிப்
பார்க்கிறார்கள்.

%%K t03-c2-11-1 p
11. --- முதலில், இதோ சில \VUgras{சிறுவர்கள்} \textsuperscript{(35)} ; அவர்கள்
\VUgras{இசைக்கலைஞர்களுக்கு} \textsuperscript{(36)} முன்னால் செல்கிறார்கள்.
பிறகு \VUgras{மேயர்} \textsuperscript{(37)}, துணை மேயர்கள், சிறு நகரத்தின்
\VUgras{ஊராட்சி மன்றம்} வருகிறார்கள். கடைசியாக \VUgras{தீயணைப்பு வீரர்கள்}
\textsuperscript{(38)} நடக்கிறார்கள்.

%%K t03-tit-12 sub
\VUsaut{5.0mm}
\VUpk{10.2pt}{40}{திருவிழா.}
\VUsaut{3.6mm}

%%K t03-c2-12-1 p
12. --- \VUgras{திருவிழா மைதானத்தில்} \textsuperscript{(39)} ஏராளமான மக்கள்
இருக்கிறார்கள். பலவகைக் \VUgras{கொட்டகைகளைப்} பார்க்க வருகிறார்கள் :
\VUgras{மரக் குதிரைகள்} \textsuperscript{(40)} ; \VUgras{சர்க்கஸ்}
\textsuperscript{(41)}, அங்கே காட்சி ஏற்கெனவே தொடங்கி விட்டது ;
\VUgras{பொது ஆட்டம்} \textsuperscript{(42)}, அங்கே நகரத்து இளைஞர்களும்
இளம் பெண்களும் \VUgras{நடனத்தின்} இன்பத்தை அனுபவிக்கிறார்கள்.

%%K t03-c2-13-1 p
13. --- உள்பக்கத்தில் \VUgras{காளைக் களத்தைப்} \textsuperscript{(44)}
பார்க்கிறோம் ; அங்கே துணிச்சலான ஸ்பானிய \VUgras{காளைச் சண்டை வீரர்}
சினம் கொண்ட \VUgras{காளையுடன்} \textsuperscript{(45)} போரிடுகிறார். பிறகு
\VUgras{சறுக்குப் பாதை} \textsuperscript{(49)}, \VUgras{சுடு பயிற்சி இடம்}
\textsuperscript{(46)} — அங்கே \VUgras{துப்பாக்கிகளை} உபயோகிக்கவும்
\VUgras{இலக்கை} நோக்கிச் சுடவும் பயிற்சி செய்கிறார்கள்.

%%K t03-c2-14-1 p
14. --- அருகில், இதோ \VUgras{திருவிழா நாடக அரங்கு} \textsuperscript{(47)} ;
அங்கே \VUgras{நகைச்சுவை நாடகமும்} \VUgras{சோக நாடகமும்} நடிக்கிறார்கள்.
மிக அருகில் \VUgras{மிகப் பெரிய மனிதரின்} \textsuperscript{(48)} கொட்டகையும்
\VUgras{மிகச் சிறிய மனிதரின்} கொட்டகையும் இருக்கின்றன. முதலாமவரின்
\VUgras{உயரம்} இரண்டு மீட்டர் பதினைந்து சென்டிமீட்டர் ; இரண்டாமவர்
ஒரு குழந்தையின் அளவுகூட இல்லை.

%%K t03-c2-15-1 p
15. --- கடைசியாக, இதோ பெரிய \VUgras{மிருகக் காட்சிச்சாலை}
\textsuperscript{(50)} ; அதில் \VUgras{காட்டு விலங்குகள்} இருக்கின்றன —
வலிமையான \VUgras{நகங்களை} உடைய \VUgras{புலி} \textsuperscript{(51)},
அடர்ந்த \VUgras{பிடரி மயிரை} உடைய \VUgras{சிங்கம்} \textsuperscript{(52)},
இரண்டு \VUgras{ஒட்டகச்சிவிங்கிகள்} \textsuperscript{(53)}, தன்
\VUgras{தும்பிக்கையை} மிகவும் திறமையாக உபயோகிக்கும் \VUgras{யானை}
\textsuperscript{(54)}.

%%K t03-c2-16-1 p
16. --- அந்த விலங்குகளைப் பழக்கும் \VUgras{விலங்கு பழக்குபவர்}
\textsuperscript{(55)} கொட்டகையின் கதவருகில் இருக்கிறார்.

\VUsaut{6.0mm}
\VUfilet{22mm}
